% Insert for resume3An26w.tex; the parent supplies headers and page numbers.
% Build: platex -interaction=nonstopmode -halt-on-error smalltest2_review.tex
% Then: dvipdfmx smalltest2_review.dvi
\documentclass[a4paper,10pt,dvipdfmx]{jsarticle}
\usepackage{../an26handout}

\fancyhead[L]{\small\color{ink}解析学 II・小テスト}
\fancyfoot[C]{\small\color{ink}\thepage}

\renewcommand{\handout}[2]{%
  \def\handoutnumber{第#1回}%
  \begin{center}
    {\Large\normalfont\mdseries\color{ink}第#1回\quad #2}\par
    \vspace{2pt}{\color{accent}\rule{0.88\linewidth}{1.2pt}}
  \end{center}\vspace{2pt}}
\renewcommand{\topic}[1]{\par\vspace{5pt}%
  {\large\normalfont\mdseries\color{accent}#1}\par\vspace{1pt}}

\begin{document}
\pagestyle{empty}
\handout{2}{小テスト：有理関数の積分}

学籍番号：\underline{\hspace{40mm}}\hfill
氏名：\underline{\hspace{40mm}}

次の不定積分を求めてください。計算過程も書き、積分定数を付けてください。
問1は30点、問2は40点、問3は30点、合計100点です。

\topic{問題}
\begin{enumerate}[label=\color{accent}\arabic*.,leftmargin=2.4em,itemsep=9pt]
\item $\displaystyle\int\frac{1}{x^2-3x+2}\,dx$\hfill（30点）

\item $\displaystyle\int\frac{x}{x^2-2x+1}\,dx$\hfill（40点）

\item $\displaystyle\int\frac{1}{x^2-2x+2}\,dx$\hfill（30点）
\end{enumerate}

\par\vspace{4mm}
{\color{accent}\rule{\linewidth}{0.6pt}}
\topic{小テストの解答}
以下、$C$ は区間ごとに選ぶ積分定数、$\log$ は自然対数を表します。

\begin{enumerate}[label=\color{accent}\arabic*.,leftmargin=2.4em,itemsep=9pt]
\item 分母を因数分解し、部分分数に分けます。
$$
 x^2-3x+2=(x-1)(x-2),\qquad
 \frac1{(x-1)(x-2)}=-\frac1{x-1}+\frac1{x-2}.
$$
したがって、
$$
 \int\frac{1}{x^2-3x+2}\,dx
 =-\log|x-1|+\log|x-2|+C
 =\log\left|\frac{x-2}{x-1}\right|+C.
$$
原始関数を考える区間は $(-\infty,1)$、$(1,2)$、$(2,\infty)$ です。

\item 分母は $(x-1)^2$ です。分子を $x=(x-1)+1$ と書くと、
$$
 \frac{x}{x^2-2x+1}
 =\frac{(x-1)+1}{(x-1)^2}
 =\frac1{x-1}+\frac1{(x-1)^2}.
$$
したがって、
$$
 \int\frac{x}{x^2-2x+1}\,dx
 =\log|x-1|-\frac1{x-1}+C.
$$
原始関数を考える区間は $(-\infty,1)$、$(1,\infty)$ です。

\item 分母を平方完成し、$u=x-1$ と置きます。
$$
 x^2-2x+2=(x-1)^2+1,\qquad du=dx.
$$
したがって、
$$
 \int\frac{1}{x^2-2x+2}\,dx
 =\int\frac{du}{1+u^2}
 =\arctan u+C
 =\arctan(x-1)+C.
$$
分母はすべての実数 $x$ で正なので、実数全体で原始関数を考えられます。
\end{enumerate}
\end{document}
